\chapter{OntoUML Overview}

\section{Motivation}

Conceptual models are used to describe domains of interest, support communication among stakeholders, and provide a basis for analysis and system development. Their usefulness thus depends on how meaning can be conveyed through the use of a language's modeling constructs. General-purpose modeling languages provide a broad set of constructs, yet their semantics may leave important ontological commitments implicit. As a consequence, a model can be syntactically acceptable while obscuring distinctions that are essential to the domain, such as the difference between an object and a relationship, or between an entity type and its instances \cite{guizzardi_wagner_2010_ufo_foundation}.

Foundational ontologies address this problem by providing general (i.e., domain-independent) categories and principles against which domain models and modeling languages can be examined. UFO was developed with this purpose for conceptual modeling. Its application to business modeling illustrates how ontological distinctions can reveal shortcomings in existing methods, including the conflation of events with activities and the treatment of commitments as if they were events \cite{guizzardi_wagner_2005_business_modeling}. These analyses motivate the development of modeling languages whose constructs have explicit and coherent interpretations grounded in an ontological theory.

\section{Ontology-Driven Conceptual Modeling}

Ontology-Driven Conceptual Modeling (ODCM) employs ontological theories as foundations for designing, evaluating, and using conceptual modeling languages. The ontological categories are reflected in the language's modeling constructs, while constraints derived from the theory restrict the models that can be formed with those constructs. This approach gives the constructs explicit semantics and provides criteria for assessing whether a language can represent a domain faithfully and whether its models can be understood by their intended users \cite{guizzardi_wagner_2010_ufo_foundation}.

ODCM therefore concerns more than the use of an ontology as a vocabulary for a particular application domain. It uses a foundational theory to guide the construction and interpretation of conceptual models across domains. The resulting models make their ontological commitments explicit, and their well-formedness conditions help prevent combinations of constructs that conflict with those commitments. UFO was developed as such a foundation, and OntoUML is an example of a language engineered according to this approach \cite{towards_ontological_foundations_for_conceptual_modeling__the_unified_foundational_ontology__ufo__story_2015}. 

\section{Relationship between OntoUML and UFO}

UFO provides the ontological theory that informs the distinctions represented in OntoUML. The language originated as an ontologically well-founded version of the UML class-diagram fragment for structural conceptual modeling. Its constructs were designed to reflect the distinctions in UFO-A, the part of UFO concerned with endurants\footnote{Endurants are object-like entities that exist in time, possibly changing while maintaining their identity. These include independent endurants (e.g., a Person, a Car) but also existentially depedendt endurants (e.g., a particular headache, a particular Marriage).}, and its syntactic constraints were intended to ensure that valid models conform to the corresponding ontological commitments \cite{towards_ontological_foundations_for_conceptual_modeling__the_unified_foundational_ontology__ufo__story_2015}.

In OntoUML, these distinctions are represented through stereotypes and constraints. For example, stereotypes distinguish \textit{kinds}, \textit{roles}, \textit{phases}, and other categories of types, while constraints govern how those categories may be related in a taxonomy. The relation is thus one of theoretical foundation and language realization: UFO provides the ontological categories and principles, while OntoUML provides modeling constructs with which to represent them in conceptual models. OntoUML does not define UFO itself; it makes selected UFO distinctions available in a UML-based language. By using the language, modelers employ the foundational distinctions to build well-founded conceptual models for their domains of interest.

OntoUML has evolved over the years through continued research and experience with its application in different domains. This development has led to refinements in the language, its underlying theory, and the methods and tools that support its use. The evolution of OntoUML and UFO reflects an ongoing interaction between theoretical work and the needs identified in practice \cite{towards_ontological_foundations_for_conceptual_modeling__the_unified_foundational_ontology__ufo__story_2015,endurant_types_in_ontology_driven_conceptual_modeling__towards_ontouml_2_0_2018}.

\section{Relationship between OntoUML and UML Class Diagrams}

OntoUML is defined as a UML profile for ontology-driven conceptual modeling. It retains the familiar class-diagram form and builds on UML modeling elements, including classes, associations, properties, and generalizations. OntoUML adds stereotypes and semantically motivated constraints that assign more specific ontological interpretations to these elements \cite{guizzardi_wagner_2010_ufo_foundation,towards_ontological_foundations_for_conceptual_modeling__the_unified_foundational_ontology__ufo__story_2015}.

A UML class diagram can represent types and their structural relationships without prescribing a foundational ontology for interpreting them. In OntoUML, by contrast, a class is assigned an ontological category through its stereotype, and constraints govern the admissible relations among categories. Thus, a class stereotyped as a kind represents an identity-providing type, whereas a class stereotyped as a role represents a type whose instances are classified contingently in a relational context. These stereotypes do not merely provide labels for diagram elements; they carry semantics grounded in UFO and affect which models are well-formed.

This relationship preserves the notation and modeling infrastructure of UML while providing an ontology-grounded interpretation for conceptual models. OntoUML class diagrams are therefore used to describe the types and relations of a domain, rather than to prescribe the design of software classes. The UML profile mechanism supplies the language extension point; UFO-based distinctions and constraints characterize the OntoUML extension.

% \section{Model Semantics}

% Discuss here that models can be interpreted using a snapshot semantics or a growing-block semantics.
% JP: commented this out as I don't think this is the place for this; to appreciate this discussion the reader will need much more background (including UFO foundation from Chapter 8)

\section{Modeling Principles}

While UML emerged from object-oriented analysis and design methods and incorporates concepts from object-oriented programming, it can also be employed as a conceptual modeling language beyond software design. 

OntoUML is positioned in this landscape as an ontology-driven conceptual modeling language that uses UML's class-diagram notation and extension mechanisms to represent domain structures, while leaving implementation-specific concerns to subsequent design and transformation activities \cite{guizzardi_wagner_2010_ufo_foundation,towards_ontological_foundations_for_conceptual_modeling__the_unified_foundational_ontology__ufo__story_2015}.

Conceptual modeling and object-oriented (OO) design modeling address related but distinct concerns. Conceptual models describe the structure of a domain and the distinctions considered relevant for understanding it. OO design models, in contrast, describe software structures and are shaped by the constructs and restrictions of a particular implementation environment. Consequently, the adequacy of a conceptual model should be assessed in relation to the domain it represents, rather than by how directly its constructs correspond to those available in an object-oriented programming language or a database schema \cite{guizzardi_wagner_2010_ufo_foundation,guidoni_almeida_guizzardi_2020_relational_schemas}.

This distinction is relevant to the use of abstraction mechanisms in UML (and leveraged in OntoUML) such as multiple inheritance, dynamic classification and orthogonal generalization sets. These mechanisms allow a conceptual model to represent different criteria of classification that apply to the same domain entities. For example, persons may be classified (dynamically) by life phase and by citizenship, so that a person may instantiate both Adult and Brazilian Citizen (at a particular point in time). Such classifications may overlap, and an individual may change from one classification to another without ceasing to be the same individual. These structures support modularity and reuse in conceptual models and should not be excluded merely because a particular programming language or implementation strategy offers limited support for them \cite{guidoni_almeida_guizzardi_2020_relational_schemas}. Transformations to relational schemas, for example, must account for features such as multiple inheritance, overlapping and incomplete classifications, orthogonal hierarchies, and dynamic classification. The conceptual model should preserve these domain distinctions; the choice of realization strategy should address how they are represented in the target technology \cite{guidoni_almeida_guizzardi_2020_relational_schemas}.

OntoUML can be understood as a pattern language because its modeling constructs do not occur freely, but are combined in recurring configurations whose structure and constraints express ontological micro-theories. These ontology-based patterns provide reusable solutions to recurring conceptual modeling tasks, while inductive modeling rules can guide their construction through a sequence of modeling decisions, including in tool-supported dialogues with modelers. A pattern language organizes such patterns and their relations to help modelers select and apply them; anti-patterns, in turn, describe recurrent structures that may lead to unintended interpretations. The approach also supports model stability: patterns for orthogonal classifications, multi-level structures (with powertypes), reified aspects, and non-sortal types help accommodate changes in requirements and domain conceptualization. In this way, OntoUML patterns support reuse, guide model construction, and help manage complexity while retaining an explicit ontological basis \cite{guizzardi_gracas_guizzardi_2011_ontouml_patterns,guizzardi_2014_ontological_patterns,guizzardi_almeida_2020_stability_patterns,carvalho_almeida_guizzardi_2016_powertype_pattern}.